\clearpage\thispagestyle{empty}
\pdfbookmark[0]{Quellen zur Besprechung von Knitlmayers Schrift über Planetenabstände}{knitlmayer1810-sources}
\begin{center}\textbf{Quellen zur Besprechung von Knitlmayers Schrift über Planetenabstände}\end{center}
\noindent Diese KI-übersetzte Fassung in modernem Deutsch folgt Carl Friedrich Gauß, \textit{Werke}, Band~VI (Göttingen, 1874), S.~528--530. Die Besprechung erschien zuerst ohne Verfassernamen unter Brünn in den \textit{Göttingischen gelehrten Anzeigen}, Stück 179, 10. November 1810, S.~1777--1782.
\par\smallskip\noindent Gauß bespricht Ferdinand Knitlmayers \textit{Versuch einer genauen Darstellung des Progressionsverhältnisses der Planeten- und Trabantenabstände von ihren Centralkörpern.} Constant von Wurzbachs \textit{Biographisches Lexikon des Kaiserthums Oesterreich}, Band~12 (1864), S.~153, nennt Gastl, Brünn, 1808 als Verlagsangabe.
\par\smallskip\noindent In der Werkausgabe steht \emph{geocentrische Progression}, im Erstdruck \emph{geometrische Progression}. Die Übertragung folgt der Werkausgabe. In beiden Drucken steht vor dem Bruch in der Kritik an der Dichteberechnung ein Minuszeichen. In der Formel zu den Trabanten fehlt bei der Kubikwurzel in beiden Drucken der Strich über dem Radikanden.
\par\smallskip\noindent \href{https://gdz.sub.uni-goettingen.de/id/PPN236007467}{Text der Werkausgabe}. Erstdruck, \href{https://images.p.sub.uni-goettingen.de/iiif/image/gdz:PPN319721507_1810:00001811/full/1650,/0/default.jpg}{S.~1777}, \href{https://images.p.sub.uni-goettingen.de/iiif/image/gdz:PPN319721507_1810:00001814/full/1650,/0/default.jpg}{S.~1780}, \href{https://images.p.sub.uni-goettingen.de/iiif/image/gdz:PPN319721507_1810:00001815/full/1650,/0/default.jpg}{S.~1781}. \href{https://repository.uibk.ac.at/filestore/servlet/GetFile?id=CHFZTKXMVZVIMCGAEKCS&convert=image/jpg&scale=3.33}{Wurzbachs Eintrag}.
\par\smallskip\noindent Übersetzung und Satz mit KI erstellt.
\clearpage
